% Drop-in section. Requires natbib or equivalent citation commands.
% The final paragraph assumes only the stated full-image, untrailed-star capability.
\section{Previous work}
\label{sec:previous-work}

Astrometric analysis of a hemispheric fisheye image involves both identifying
catalogue stars and determining the mapping between image coordinates and sky
directions. These tasks should be distinguished when comparing methods: a
calibration valid across the complete image may be obtained from star
identifications established in smaller regions or supplied during an initial
calibration. Blind plate solving, exemplified by Astrometry.net
\citep{Lang2010}, establishes image--catalogue correspondences through geometric
pattern matching and statistical verification. Whether this identification step
operates on the complete fisheye field, and what information is required to
initialise it, are therefore central to the comparison.

Several all-sky pipelines use conventional plate solvers to establish local
correspondences before deriving a global calibration. \citet{Hung2024} apply
Astrometry.net to five selected $280\times280$-pixel subimages of their fisheye
night-sky images. \citet{Yin2025} reproject regions of an all-sky image into
views suitable for Astrometry.net, combine the resulting stellar
identifications, and construct a global coordinate transformation. A related
but distinct strategy is used in the FRIPON calibration of
\citet{Jeanne2019}: SExtractor and SCAMP provide initial stellar
identifications at elevations above $30^\circ$, where distortion is less severe,
after which a fisheye model is fitted using stellar measurements accumulated
across calibration exposures. Thus, a final transformation covering the full
sky does not in itself imply that the initial plate solution was obtained from
the entire hemispheric image in one operation.

Other approaches perform global calibration with instrument or observational
information available at initialisation. \citet{Barghini2019} describe an
automated all-sky calibration in which a simplified projection model generates
predicted catalogue-star positions, which are correlated with detected sources
and progressively refined. Their initialisation uses an approximate plate
scale, a prescribed image orientation, and the observing time and station
coordinates. The ORION software of \citet{AntunaSanchez2022}
likewise derives full-image angular-coordinate maps from stellar positions;
its automatic star search uses a previous calibration or an approximate initial
mapping. \citet{Yang2025} fit a global distortion model for super-wide-field
images, beginning with 50 manually identified stars and expanding the matched
set iteratively. These studies demonstrate that neither global fisheye
calibration nor automated refinement of a whole-image model is, by itself,
a new capability.

Related work in celestial navigation addresses star identification in very
wide fields with external assistance. \citet{Li2022} identify stars in
$180^\circ$ fisheye imagery using precise point positioning and inertial
navigation information to predict their image locations. The problem is thereby
constrained by external estimates of position and attitude. Calibration from
stellar motion provides another route: \citet{TeagueChahl2024} bootstrap camera
geometry from two time-separated images, although their experimental cameras
have a horizontal field of view of $53.5^\circ$. In the Auto-Cal preprint,
\citet{Kapali2025} combine laboratory calibration of fisheye intrinsics with
stellar tracks measured through image sequences to determine and update camera
orientation. These temporal and instrument-assisted methods address related
calibration problems but require inputs beyond an individual stellar image.

The present work addresses plate solving of an entire $180^\circ$ fisheye image
containing untrailed stellar sources, without dividing the image into
independently solved patches. This distinguishes the processing strategy from
the explicitly subdivided approaches described above, while the global methods
of \citet{Barghini2019}, \citet{AntunaSanchez2022}, and \citet{Yang2025} remain
relevant points of comparison. The comparison therefore concerns star identification, initialisation
requirements, and performance across the hemispheric field.
